\documentclass[a4paper,10pt]{article}
\usepackage[utf8]{inputenc}

%Preámbulos de las figuras

%--------------------------------------------
\usepackage{polynom}
%--------------------------------------------
\usepackage{nicematrix,tikz}
\usetikzlibrary{fit}
%--------------------------------------------
\usepackage[table, svgnames]{xcolor}
\usepackage{hhline, booktabs}
%--------------------------------------------
\usepackage{amsmath,gensymb}
\usetikzlibrary{matrix,positioning,calc}
\definecolor{COLOR1}{HTML}{6B99B9}
\definecolor{COLOR2}{HTML}{DD9CC4}
\tikzset{
    M/.style={
        matrix of math nodes,
        nodes={draw=none,anchor=center,minimum size=1.75em},
        nodes in empty cells,
    }
}
%Drawing code nesting in definition with parameters
\def\Det[color #1 size #2 line #3]#4#5{
    \draw[line width=#3,draw=#1, opacity=0.7]
    let \p1 = ($(#5)-(#4)$), % To access cartesian coordinates x, and y.
        \n2 = {atan2(\x1,\y1)}, % to get the angle of the line betwen #5 to #4
    in
    (#5.center)
        ++(90-\n2+90:#2/2)
        arc (90-\n2+90:90-\n2+90-180:#2/2)
        -- ($(#4.center)+(90-\n2-90:#2/2)$)
        arc (90-\n2-90:90-\n2-90-180:#2/2)
        -- cycle;
}
%--------------------------------------------
\usepackage{amssymb}
\newcommand{\ff}{\mathtt{f}}
%--------------------------------------------
\newcolumntype{A}{>{\columncolor{red!20}}c}
\newcolumntype{B}{>{\columncolor{blue!20}}c}

%%------------------------------------------------------------------------


\pagestyle{empty}
\begin{document}
\setlength\extrarowheight{2pt}
\setlength{\arrayrulewidth}{0.6pt}
\arrayrulecolor{Gainsboro!50!Lavender}
\begin{tabular}{>{\columncolor{Gainsboro!50!Lavender}}l|!{\color{RoyalBlue}\vrule width 0.6pt}c|c|c!{\color{RoyalBlue}\vrule width 0.6pt}}
\rowcolor{Gainsboro!50!Lavender} \multicolumn{1}{c!{\hspace{\dimexpr\doublerulesep + \arrayrulewidth}}}{\cellcolor{PapayaWhip}} & \multicolumn{1}{!{\hspace{-\arrayrulewidth}}c}{headline1} & \multicolumn{1}{c}{headline2} &\multicolumn{1}{c}{ headline3} \\\addlinespace[0.5ex]
\hhline{-|>{\arrayrulecolor{RoyalBlue}}|---|}
row1 & a & b & c \\
\hhline{>{\arrayrulecolor{Gainsboro!50!Lavender}}-|>{\arrayrulecolor{RoyalBlue}}|---|}
  row2 & d & e & f \\
\hhline{>{\arrayrulecolor{Gainsboro!50!Lavender}}-|>{\arrayrulecolor{RoyalBlue}}|---|}
 row3 & g & h & i \\
\hhline{>{\arrayrulecolor{Gainsboro!50!Lavender}}-|>{\arrayrulecolor{RoyalBlue}}|---|}
\end{tabular}
\end{document}